\begin{abstract}

Radio and X-ray observations track the response of intracluster gas to merging events, but the merger-history information retained in realistic observing conditions remains uncertain.
We train convolutional neural networks on synthetic images of 352 systems in TNG-Cluster to infer \emph{pseudo-TSC}, an indicator of recent merger activity derived from the merger score.
This proxy strongly correlates with catalogue collision times at $z=0$ ($\rho=0.96$).
Idealized observations combining three orthogonal projections yield mean $R^2$ with calculated pseudo-TSC values of $R^2=0.485$, 0.511, and 0.545 for 144 MHz radio, X-ray, and joint inputs, respectively.
Radio predictions remain informative within mass bins and after normalization of integrated emission, indicating information content beyond total brightness.
We also find that dense gas, including a significant cold component, dominates the unfiltered synthetic radio images, motivating a density cut to restrict emission to diffuse gas.
We construct approximate LOFAR and Chandra forward models, and apply them to simulated cluster images at both redshifts, creating a training dataset closer to real-world observations.
For individual simulated observations, five-model ensembles reach $R^2=0.203$ for radio images of physical size 3 Mpc, and $R^2=0.406$ for X-ray images.
Adding 1 Mpc radio images to X-rays yields no statistically-significant improvement.
Applying the CNN to the Local Volume-Complete Cluster Survey (LoVoCCS) yields radio predictions for 26 clusters, and overlapping X-ray predictions for 15 clusters.
For the overlapping sample, the merger time rankings between radio and X-ray correlate ($\rho=0.61$), but radio assigns systematically older proxy ages, producing a median age offset of 0.44 Gyr.
These results indicate the presence of recoverable information about recent merger activity in simulated cluster observations, and show that X-rays are a stronger predictor than radio under the tested observing conditions.
However, independent calibration is needed to establish accuracy of the predicted physical collision ages.

\end{abstract}
